% META: {"id": "TEXP-GRA-EX-018", "chapitre": "TEXP-GRAPHES", "type_objet": "exercice", "capacites": ["TEXP-GRA-C3"], "capacites_codes": ["C3"], "methodes": ["M2","M5"], "parcours": 2, "competences": ["modeliser","raisonner"], "duree_min": 25, "mode_creation": "ex_nihilo", "sources_inspiration": [], "fichier_tex": "chapitres/TEXP-GRAPHES/exercices/TEXP-GRA-EX-018.tex", "corrige_tex": "chapitres/TEXP-GRAPHES/corriges/TEXP-GRA-CO-018.tex", "status": "approved"}
% BEGIN-VERIFY
% from sympy import *
% x = symbols('x')
% # boite de base carree 20x20, hauteur x, volume V = x(20-2x)^2 sur [0,10]
% V = x*(20 - 2*x)**2
% Vp = diff(V, x)
% crit = solve(Vp, x)
% assert Rational(10,3) in crit
% Vmax = V.subs(x, Rational(10,3))
% assert Vmax == Rational(16000, 27)
% assert V.subs(x, 0) == 0
% assert V.subs(x, 10) == 0
% END-VERIFY
\begin{exercice}{TEXP-GRA-EX-018}{3}{25}
On decoupe aux quatre coins d'une plaque carree de cote $20$ cm des carres de cote $x$ cm, puis on replie pour former une boite ouverte.
\begin{enumerate}
  \item Exprimer le volume $V(x)$ de la boite en fonction de $x$.
  \item Pour quelle valeur de $x$ le volume est-il maximal ? Quel est ce volume maximal ?
\end{enumerate}

\end{exercice}
